\section{绪论}

\subsection{问题的提出：影像清楚之后，临床还需要什么}

MRI能用多种组织对比把人体内部结构显示得很清楚，但“看得见病灶”，离完成临床任务还有一段距离。影像科医生报出一处异常信号，临床接着要问的往往是：这究竟是不是病变，侵袭性有多大，治疗有没有效，以后会不会复发。小体积急性脑梗死就是个例子。单看血管改变未必能认出梗死灶，高燕等因此把弥散加权成像结果当作评价其他检查的重要参照\cite{gao2025infarct}。肿瘤也一样，形态和大小很难单独回答病理分级，周定燕等于是把扩散与动态增强参数放在一起，去预测前列腺癌的Ki-67表达和Gleason评分\cite{zhou2024prostate}。两个例子指向同一件事：临床缺的不是更细的解剖图，而是能描述组织功能、微结构和疾病风险的证据。

MRI的技术演进刚好提供了这种可能。扩散加权成像（diffusion-weighted imaging，DWI）把水分子运动变成图像对比，动态对比增强磁共振成像（dynamic contrast-enhanced MRI，DCE-MRI）把对比剂在组织中的动态分布变成灌注与通透性信息，功能磁共振成像（functional MRI，fMRI）则借血氧水平依赖信号间接观察脑活动。三者都没有取代常规T1WI与T2WI，只是在其形态信息之外另开了观察维度。要理解它们，光记住缩写和临床用途不够，还得回到信号怎么产生、梯度怎么编码、图像怎么重建这一层\cite{wu2024mri}。

\subsection{本文的思路与基本判断}

刚接触这些技术时，我很容易把它们当成设备菜单里不断增加的“高级功能”：参数越多、图像越花，就觉得技术越先进。读了一批原理教材和临床文献之后，这个印象慢慢松动了。真正值得追问的其实是三个前后相扣的问题——某种物理或生理过程怎么被编码成信号，信号怎么经模型变成参数，参数又怎么在一个具体病例上改变了诊断或治疗判断。本文因此不按技术清单排列，而是沿“信号—参数—证据—决策”这条链往下走。

文章先交代射频激发、弛豫、梯度空间编码与K空间重建这些共同基础，再分别讨论扩散成像、DCE-MRI和BOLD-fMRI的原理、改进脉络、典型案例与局限，最后看多参数成像、影像组学和深度学习落在决策环节的什么位置。贯穿全篇的判断只有一句：MRI从“看见病灶”走到“支持决策”，关键不在于图像数量或模型复杂度，而在于新增的信息能否回答一个明确的临床问题，并且这种回答经得起重复与验证。

这条线索也让我重新理解本专业的交叉属性。物理原理决定能测到什么，工程方法决定测得稳不稳，医学问题决定测出来的东西有没有用。三者缺一环，技术就容易停在“能出图”的阶段。
